\documentstyle[%


righttag%


,11pt%


,amssymb%


,russian%               remove in Eng.


,izvru%                 Set izven% in Eng.


% ,izvru-ks             for brief communications


]{amsart}


% File from author


% Math definitions





% \newcommand{\ov}{\overline}


\newcommand{\End}{\operatorname{End}}


\newcommand{\vraisup}{\operatorname{vrai~sup}}


\newcommand{\tts}[1]{}





\theoremstyle{definition}


\theorembodyfont{\rm}


\newtheorem{examnonum}{\hskip\parindent Пример}


\renewcommand{\theexamnonum}{}





%\hoffset=-15mm%                  For Eng.





\setcounter{page}{60}


\god{1998}


\nomer{2~(429)} %                Remove in Eng.


%\nomer{Vol.\,42, No.\,2}%        Set in Eng.


% \str{75--81}%                    Set in Eng.


\udk{512.55}


\author{\it И.И.\,Сахаев, Х.И.\,Хакми}


% NB:   ^^ remove in Eng.


\title{О сильно регулярных модулях и кольцах}





\thanks{Работа выполнена при финансовой поддержке Российского фонда


фундаментальных исследований (грант \No.\,96-01-00717).}


\date{}


% \pagestyle{myheadings}%         Set in Eng.





\pagestyle{plain} %              Remove in Eng.


\begin{document}





% \thispagestyle{plain}%          Set in Eng.





\maketitle





Работа посвящена изучению строения сильно регулярных модулей


и колец, в частности, сильно регулярных проективных и инъективных


модулей. Тем самым существенно дополняются исследования,


проведенные в [1]--[4]. Рассматриваются ассоциативные кольца с


единицей и унитарные модули над ними. Будем обозначать через


$J(R)$ и $J(M)$ радикал Джекобсона соответственно кольца


$R$ и модуля $M$.





Напомним определения из [1]--[4].


Правый $R$-модуль $M$ называется слабо регулярным (соответственно


сильно регулярным), если для каждого подмодуля $N$


модуля $M$, $N \nsubseteq J(M)$, существует подмодуль


$Q$ модуля $N$ такой, что $Q \not\subset J(M)$ и $M=Q \oplus Q'$


(соответственно $M=N \oplus Q'$). Кольцо $R$ называется


слабо регулярным (сильно регулярным), если оно, рассматриваемое


как правый $R$ модуль, является слабо регулярным (сильно


регулярным). Правый $R$-модуль $M$ называется $n$-регулярным,


если каждый его подмодуль $N$, порожденный $n$ элементами


и $N \not\subset J(M)$, является прямым слагаемым модуля $M$.


Правый $R$-модуль $M$ называется локальным, если радикал


$J(M)$ является наибольшим подмодулем модуля $M$, и полупростым,


если он является прямой суммой простых правых $R$-модулей.


Если $M=J(M)$, то правый $R$-модуль $M$ называется


радикальным. Радикальный правый $R$-модуль $M$, очевидно,


является сильно регулярным, а регулярное (в смысле Неймана)


кольцо является $n$-регулярным для любого целого $n$.


Переходим к изложению результатов.





\begin{thm} Если правый $R$-модуль $M$


$2$-регулярен, $M \ne 0$, то либо $M=J(M)$, либо $J(M)=0$, либо


$M$ локальный{\em ;} и наоборот, если $M$ есть $1$-регулярный правый


$R$-модуль, удовлетворяющий одному из предыдущих ограничений на


модуль $M$, то $M$ есть $2$-регулярный правый $R$-модуль.


\end{thm}





{\bf Доказательство} проведем в нескольких пунктах.





1. Пусть $M \ne J(М)$, $J(M) \ne 0$. Покажем, что в этом


случае $M$ является локальным правым $R$-модулем. Предположим,


что $M$ --- не локальный $R$-модуль.


Так как $M \ne J(M)$, то существует такой $a \in M$, что


$aR \not\subset J(М)$.





2.1. Рассмотрим случай $M=aR$. Поскольку


$aR$ --- не локальный правый $R$-модуль, то существует $bR \subset aR$,


$bR \not= aR$, $bR \not\subset J(aR)$. Покажем, что $J(aR)=J(bR)$.


Пусть существует такой $z \in J(aR)$, что $bR+zR \ne bR$.


Ввиду 2-регулярности $R$-модуля $aR$ имеем $aR=( bR+ zR ) \oplus N$.


В силу следствия 9.1.3~(а) [5] подмодуль $zR$ косущественен


в $aR$, поэтому $aR=bR \oplus N$. Отсюда $bR+zR =(bR+zR) \bigcap


aR=(bR+zR) \bigcap (bR \oplus N)=bR$, пришли в противоречие


с предположением $bR+zR \ne bR$. Значит, $z \in J(aR)$ влечет


$z \in bR$, поэтому $J(aR)=J(bR)$ и $aR=bR \oplus N$. Если $N=0$,


то получим $aR=bR$ и приходим в противоречие


с предположением $aR \ne bR$.





2.2. Пусть $N \ne 0$. Так как $J(aR)=J(M) \ne 0$, то


существует $z \ne 0$, $z \in J(aR)$, $q \in N$, $q \ne 0$,


$zR \oplus qR \not\subset J(aR)$. Снова в силу 2-регулярности


модуля $aR$ вытекает, что $aR= zR \oplus qR \oplus N^*=


N_1 \oplus zR \oplus qR \oplus N_2$, где


$N_1 \oplus zR =bR$, $qR \oplus N_2=N$. Отсюда в силу


косущественности подмодуля $zR$ получим $aR= N_1 \oplus


qR \oplus N_2$, поэтому $z=0$ и $J(aR)=0$, пришли к противоречию,


ибо $J(aR) \ne 0$. Следовательно, $M=aR$ --- локальный


правый $R$-модуль.





2.3. Рассмотрим общий случай: $a \in M$, $aR \not\subset J(M)$,


существует $z \in J(M)$, $zR \not\subset aR$, $aR+zR \ne aR$,


$aR+zR \not\subset J(M)$. Тогда, как и в п.~2.1, в силу


2-регулярности модуля $M$ имеем $M= (aR+zR) \oplus Q =


aR \oplus Q, aR+zR=aR$, пришли к противоречию. Значит, $J(M)=


J(aR)$ и в силу следствия 9.1.5 [5] вытекает $J(Q)=0$.


Если $Q=0$, то приходим к рассмотренному случаю п.~2.1, где


$M=aR$.





2.4. Пусть $M=aR \oplus Q$, $Q\ne 0$. Проведя рассуждения,


как в п.~2.2, приходим к противоречию, а именно, $J(M)=0$.


Следовательно, предположение, что 2-регулярный


правый $R$-модуль $M$ при условии $M \ne J(M)$,


$J(M) \ne 0$, не является локальным, неверно.





Пусть теперь $M$ является 1-регулярным правым $R$-модулем. Тогда при


выполнении одного из условий: $M=J(M)$, $J(M)=0$ или $M$ локальный ---


легко устанавливается 2-регулярность модуля $M$.


Действительно,


если $M=J(M)$, $J(M)=0$, то это очевидно. Если же $M$ --- локальный


модуль и $a_{1},a_{2} \in M$, тогда подмодуль $a_{1}R+a_{2}R$,


не содержащийся в $J(M)$, равен $M$, ибо $M$ --- локальный


$R$-модуль. $\quad\square$





\begin{thm}


Правый $R$-модуль $M$, $M \ne 0$, $M \ne J(M)$,


является сильно регулярным тогда и только тогда, когда он либо


полупростой, либо локальный $R$-модуль.


\end{thm}





\begin{pf} Прежде всего покажем, что фактор модуль $\ov M= M/J(M)$


является полупростым модулем. Пусть $\ov U \ne 0$ есть произвольный


подмодуль модуля $\ov M$.


Обозначим через $N$ такой подмодуль модуля $M$, что $J(M)


\subset N$, $N/J(M)= \ov U$. Поскольку $M$ --- сильно регулярный


$R$-модуль и $\ov U \ne 0$, то $N \not\subset J(M)$ и $M=N


\oplus Q$. Согласно следствию 9.1.5 ([5], с.\,213)


имеем $J(M)=J(N) \oplus J(Q)$ и


поскольку $J(M) \subset N$, $N \bigcap Q =0$, то $J(Q) \subseteq


J(M) \bigcap Q \subseteq N \bigcap Q$. Значит, $J(Q)=0$ и


$J(M)=J(N)$. Тогда $\ov M= N/J(N) \oplus Q/J(Q)=\ov U \oplus


\ov Q$. Следовательно, каждый ненулевой подмодуль $\ov U$


модуля $M$ выделяется прямым слагаемым в модуле $M$ и согласно


теореме 4.2 ([6]) $\overline M$ будет полупростым правым $R$-модулем.





Поскольку по условию теоремы $R$-модуль $M$ является сильно


регулярным, то он является 2-регулярным модулем и согласно


теореме 1\; $M$ --- либо локальный $R$-модуль, либо $J(M)=0$.


Если $J(M)=0$, то согласно выше доказанному $M$ --- полупростой


модуль. Если $M$ --- локальный модуль, то необходимость условий


теоремы доказана.








Пусть $M$ является полупростым $R$-модулем,


тогда согласно теореме 4.2 [6]


$M$ --- сильно регулярный правый $R$-модуль.





Пусть теперь $M$ --- локальный правый $R$-модуль.


Тогда по определению локального модуля радикал $J(M)$


является наибольшим собственным подмодулем $R$-модуля $M$ и


поэтому каждый подмодуль $N\;$ $R$-модуля $M$, не содержащийся


в радикале $J(M)$, совпадает с модулем $M$.


\end{pf}





\begin{cor} Проективный правый $R$-модуль $P$ является сильно


регулярным тогда и только тогда, когда он является либо


полупростым, либо локальным $R$-модулем.


\end{cor}





В частности, при $P=R$ справедливо





\begin{cor} Кольцо $R$ является сильно регулярным тогда


и только тогда, когда оно --- либо артиново полупростое, либо


локальное кольцо.


\end{cor}





\begin{thm} Пусть $P$ --- конечно порожденный проективный правый


$R$-модуль. Тогда следующие условия эквивалентны{\em :}


\begin{itemize}


\item[(1)] $R$-модуль $P$ сильно регулярен,


\item[(2)] кольцо $R$-эндоморфизмов $S=\End_R(P)$ правого


$R$-модуля $P$ является сильно регулярным кольцом.


\end{itemize}


\end{thm}





{\bf Доказательство.} Согласно следствию 1\; $R$-модуль $P$


является либо полупростым, либо локальным.


Пусть $P$ является локальным $R$-модулем, тогда согласно


теореме 2.5 [7] кольцо $S=\End_R(P)$ будет локальным.


Пусть теперь $P$ --- полупростой правый $R$-модуль, тогда в силу


упражнения 1 ([5], c.\,103) кольцо $S$ является регулярным (в


смысле Неймана) кольцом и поэтому $J(S)=0$.





Положим $P= P_1\oplus P_2 \oplus \dotsb \oplus P_m$,


где $P_{i}$ --- прямая сумма


всех простых прямых слагаемых $R$-модуля $P$, изоморфных


одному простому подмодулю. Согласно следствию с


и следствию b ([8], c.\,68) кольцо $S$


является


артиновым, полупростым кольцом и поэтому $S$ --- сильно


регулярное кольцо. Теперь докажем импликацию





$(2) \Longrightarrow (1)$. Пусть $S=\End_{R}(P)$


--- артиново полупростое кольцо. Тогда оно --- слабо регулярное


кольцо, и поэтому $R$-модуль $P$ будет слабо регулярным.


Пусть $e_1,\dotsc,e_m$ --- множество всех различных,


примитивных, попарно ортогональных


идемпотентов кольца $S$ таких,


что $\sum\limits_{i=1}^{m}e_{i}=1$.


Положим $U_i=e_i(P)$ $(i=1,\dotsc,m)$,


тогда, очевидно, $P=U_1 \oplus


\dotsb \oplus U_m$. Поскольку правый $R$-модуль $P$ слабо


регулярен, то в силу косущественности радикала $J(P)$


$R$-модуль $U_i$ будет сильно регулярным. Если $U_i$ ---


не простой правый $R$-модуль, то он


разложим в прямую сумму ненулевых подмодулей,


поэтому идемпотент $e_i$ разложим и непримитивен. Пришли


к противоречию. Следовательно, $U_i$ является простым


правым $R$-модулем и $R$-модуль $P$ является прямой суммой


простых правых $R$-модулей $U_i$. Поэтому $P$ --- сильно регулярный


правый $R$-модуль.





Пусть теперь $S=\End_{R}(P)$ является локальным кольцом,


тогда согласно теореме 2.5 [7] $R$-модуль $P$ является локальным


правым $R$-модулем. Импликация $(2) \Longrightarrow (1)$, а следовательно,


и теорема доказаны.





\begin{thm} Пусть $R$ --- кольцо, $I(R)$ --- инъективная оболочка


$R$ как правого $R$-модуля. Тогда следующие условия


эквивалентны:


\begin{itemize}


\item[(1)] кольцо $R$ содержится в радикале $J(I(R))$,


\item[(2)] каждый инъективный правый $R$-модуль является


радикальным.


\end{itemize}


\end{thm}





\begin{pf} $(1) \Longrightarrow (2)$. Пусть $Q$ ---


произвольный инъективный правый $R$-модуль, $q \in


Q$. Определим $R$-гомоморфизм


$\varphi_q:R\longrightarrow Q$, полагая


$\varphi (1)=q$. Рассмотрим точную последовательность


правых $R$-модулей


$$


0 \longrightarrow R \longrightarrow^\varepsilon J(I(R)),


$$


где $\varepsilon $ обозначает данное вложение $R$ в $J(I(R))$.


В силу инъективности правого $R$-модуля $Q$ существует такое


$h:J(I(R)) \longrightarrow Q$, что $ \varphi =h \cdot


\varepsilon $. Поскольку согласно теореме 9.1.4 [5] $h(J(I(R)))


\subset J(Q) $, то $\varphi (q) \in J(Q)$ для любого элемента


$q$ из $Q$. Следовательно, $Q \subseteq J(Q)$ и поэтому $Q=J(Q)$.





Импликация $(1) \Longrightarrow (2)$ очевидна, т.к. $I(R)=J(I(R))$.


\end{pf}





О том, что кольцо, удовлетворяющее условиям теоремы 4, существует,


говорит





\begin{examnonum} Пусть $Z$ --- кольцо целых чисел. Согласно ([5], гл.\,9.1,


пример 1) поле рациональных чисел $Q$ является


радикальным инъективным $Z$-модулем и инъективной оболочкой кольца~$Z$.


\end{examnonum}





\begin{cor} Если инъективная оболочка $I(R)$ кольца $R$


нерадикальна и сильно регулярна, то $R$ --- либо артиново полупростое,


либо самоинъективное локальное кольцо.


\end{cor}





\begin{pf} Поскольку по условию $I(R) \ne J(I(R))$, то


согласно теореме 4\; $R \not\subset J(I(R))$, поэтому из-за сильной


регулярности $R$-модуля $I(R)$ кольцо $R$ является прямым


слагаемым инъективного правого $R$-модуля $I(R)$. Тогда кольцо


$R$ будет самоинъективным и сильно регулярным. Согласно


следствию 2 кольцо $R$ либо локальное, либо артиново полупростое.


\end{pf}





Следующая теорема позволяет охарактеризовать артиново


полупростые кольца в терминах сильно регулярных модулей.





\begin{thm} Для кольца $R$ следующие условия эквивалентны{\em :}


\begin{itemize}


\item[(1)] $R$ --- артиново полупростое кольцо,


\item[(2)] существует сильно регулярный


свободный правый $R$-модуль ранга $2$,


\item[(3)] существует нерадикальный инъективный правый $R$-модуль,


и прямая сумма инъективной оболочки $I(R)$


кольца $R$ сильно регулярна,


\item[(4)] все инъективные и все проективные правые $R$-модули


сильно регулярны.


\end{itemize}


\end{thm}





\begin{pf} $(1) \Longrightarrow (2)$ очевидно.





$(2) \Longrightarrow (1)$. Пусть $F=\langle x_1,x_2\rangle$ ---


сильно регулярный свободный правый $R$-модуль ранга 2


c базой $x_1,x_2$. Тогда


$J(F)=x_1J(R) \oplus x_2J(R)$. Из сильной регулярности


модуля $F$ следует, что подмодуль $x_1J(R) \oplus x_2R$


выделяется прямым слагаемым в модуле $F$, т.е. $F=(x_1J(R)


\oplus x_2R) \oplus Q$. Отсюда $x_1R=x_1J(R) \oplus


N$ и $J(N)=0$, $J(x_1J(R))=J(x_1R)=x_{1}J(R)$. Поэтому


$x_1J(R)$ будет проективным $R$-модулем и


согласно теореме 9.6.3 [5] $x_1J(R)=0$. Отсюда


в силу свободности базисного элемента $x_1$ получим $J(R)=0$.


Следовательно, $R$ как сильно регулярное кольцо с нулевым радикалом


$J(R)$ будет артиновым полупростым кольцом.





$(3) \Longrightarrow (1)$. Согласно условию (3) существует


нерадикальный инъективный правый $R$-модуль $Q$. Тогда в силу


теоремы 3\; $R \not\subset J(I(R))$, поэтому согласно следствию 3


кольцо $R$ --- самоинъективное, локальное кольцо. Поскольку прямая


сумма кольца $R$ является инъективной и сильно регулярной


в силу условия (3), то, как и выше, радикал


$J(R)=0$. Поэтому $R$ --- артиново полупростое кольцо.





$(1) \Longrightarrow (3)$ и остальные импликации


очевидны.


\end{pf}





\begin{thebibliography}{9}





\bibitem{Ni}


Nickolson W.K. {\em I-rings} // Trans. Amer. Math. Soc. -- 1975.


-- V.\,207. -- P.\,361--373.


\bibitem{H1}


Хакми Х.И. {\em О некоторых свойствах I-подобных колец}.


-- 1991. -- 28~с. -- Деп. в ВИНИТИ 09.10.91, \No\,2920-B91.


\bibitem{H2}


Хакми Х.И. {\em I-подобные модули} // Изв. вузов. Математика.


-- 1993. -- \No\, 9. -- C.\,65--70.


\bibitem{H3}


Хакми Х.И. {\em Сильно регулярные и слабо регулярные кольца


и модули} // Изв. вузов. Математика. -- 1994. -- \No\,5. -- C.\,60--65.


\bibitem{Cash}


Каш Ф. {\em Модули и кольца}. -- М.: Мир, 1981. -- 368~c.


\bibitem{Car}


Картан А., Эйленберг С. {\em Гомологическая алгебра}. -- М.: Ин. лит.,


1960. -- 510~с.


\bibitem{Az}


Azumaya G. {\em F-semi-perfect modules} // J. Algebra.


-- 1991. -- V.\,136. -- P.\,73--85.


\bibitem{P}


Пирс Р.С. {\em Ассоциативные алгебры}. -- М.: Мир, 1986. -- 541~с.





\end{thebibliography}





\vskip 2cm





\postup{\it Казанский государственный}{$\qquad$\it Поступили}


\postup{\it университет}{{\it первый вариант} 25.10.1993}


\postup{\it Дамасский университет, Сирия}{{\it окончательный


вариант} 06.02.1997}





\end{document}


